\paragraph*{Confidence statements.} Each score comes with an 80\% set of possible outcomes, a statement of whether the
imaging and clinical-history scores agree, and the measured performance of the models in comparable scans. A conformal
set names the outcomes a score cannot rule out \citep{angelopoulos2023}. Write $\hat p(1) = \hat p$ for a scan's score
and $\hat p(0) = 1 - \hat p$. For each calibration scan $i$ with outcome $y_i$ the nonconformity score is
$s_i = 1 - \hat p_i(y_i)$. Because a patient contributes several scans, each calibration scan is weighted by one over
the number of that patient's calibration scans. With $\alpha = 0.2$ and $n$ calibration patients, the threshold $\hat
q$ is the smallest $s_i$ whose cumulative weight reaches $(1 - \alpha)(n + 1)$, and a new scan receives
\begin{equation}
C(\hat p) = \{\, y \in \{0, 1\} : 1 - \hat p(y) \le \hat q \,\}.
\end{equation}
As planned, each model was refitted on 70\% of the training patients and calibrated on the other 30\%, inside each
training fold on MU and once on all MU training patients for the external cohorts. The combined score's set is the
one evaluated as primary. If a new patient is exchangeable with the calibration patients, about 80\% of their sets
contain the true outcome, on average over scans but not necessarily within a subgroup
\citep{romano2020,barber2021limits}. The weighting by patient makes even the average coverage approximate, and
coverage carries over to an external cohort only if its scans are exchangeable with MU's \citep{mehrtens2023}.

% Sources: METHODS_PLAN.md M2; Supplementary Section S10.
Group-conditional (Mondrian) sets \citep{vovk2022book} compute $\hat q$ separately within each group of comparable
scans. Our groups cross the time window with the state of the enhancing tumour (absent below 0.1\,cc, without
comparison when there is no previous scan, increased when it grew by at least 1\,cc and 40\%, or not increased). Most
groups are small, so each scan is calibrated at the finest of three levels that holds at least 10 calibration
patients (its group, then its enhancement state over all windows, then all calibration scans), a fixed fallback of the
kind used for small Mondrian groups \citep{johansson2015,ding2023}. Because group calibration can raise coverage in a
group simply by making sets larger, we compared it after the planned run with marginal sets at the $\alpha$ that gives
the same share of sets holding both outcomes (exploratory).

The record carries two further descriptive statements. The agreement statement says whether the clinical-history and
imaging scores fall on the same side of their cohort medians. The reliability card, after model cards
\citep{mitchell2019}, gives for the scan's group how many scans with and without progression the models were tested
on and, with at least 10 of each (the tested tier), their AUCs, and otherwise only the counts (few cases). The card
counts the cohort's tested scans. Recomputed without the patient's own scans, it changed the tier of one MU scan.

% Sources: prism/report/facts_v3_1.py; prism/report/prompts_v3_1.py; prism/report/checker.py; prism/report/verifier.py
% (v1.1); REPORT_V3.md; REPORT_V3_1.md; METHODS_PLAN.md M3; METHODS_RESULTS.md M3; Supplementary S8; refs zeba2026,
% wiseman2017, gabriel2021, rao2025rexerr, weinberg2018, kwon2023.
